\chapter{Literature Survey}
This chapter reviews the theory and earlier work on which Folium is built. For each piece of work we note what it establishes, what it does not, and how it is used in the project.

\section{Behavioral Foundations and Observed Investor Actions}
Prospect Theory by Kahneman and Tversky \cite{b1} seeks to explain departures from the predictions of expected utility in situations of risk. Reference-dependent valuation, an asymmetric reaction to gains and losses, and non-linear probability weighting are its key features. It explains why an investor may behave as if a decision depends on the loss relative to the purchase price, even when the investment should be valued prospectively. It is a descriptive theory of choice and not a procedure for classifying a particular live situation. For Folium it gives the rationale to look at the unrealized gain or loss and the risk context as possible signals, but it does not give grounds for classifying every loss-making sale as irrational. Cumulative prospect theory \cite{b17} later added cumulative probability weighting and showed how choice under uncertainty differs from maximizing the expected value, which stresses the importance of reporting the downside tail instead of only the average simulated wealth.

In Odean's work \cite{b2}, brokerage trading records are analyzed to see whether investors realize their winning stocks more readily than their losing ones. The analysis links behavioral finance theory to actual transactions and shows that the disposition effect can be measured from trading records rather than from self-reports alone. The study does not offer a decision-time classifier, a personalized intervention or a multi-decade counterfactual simulator. Folium uses the return and the context of a position as model features; without full trading records, a single observed decision is only suggestive of the disposition effect. Shefrin and Statman \cite{b4} connect the pattern of selling winners and holding losers to mental accounting, regret aversion and self-control, and Thaler's mental accounting \cite{b18} explains how people place gains and losses in separate mental accounts. These models explain why identical wealth situations can mean different things to different investors, but they do not provide a mechanism for building a multi-class classifier that can be tested on market paths. Folium applies some of these ideas in a measurable form through portfolio features.

Bikhchandani, Hirshleifer and Welch \cite{b3} show how sequential choice can lead to informational cascades, in which people rationally adopt what others did earlier even when their own private information points the other way. Following the crowd is therefore not always irrational, since social signals carry some information. The model shows when cascades appear, but it cannot recover an individual's present intentions from market data. De Long et al. \cite{b19} show that noise traders can move prices even when rational arbitrageurs are present, so crowd measures are relevant to the market without being proof that a particular investor is herding. Folium therefore uses social indicators as proxies for herd behavior and not as conclusive evidence. Fama's review of efficient markets \cite{b20} provides the necessary balance: trading in response to information can be rational, so the system must allow rational readings instead of labelling every market-driven trade as a bias.

Barber and Odean \cite{b14} show that individual investors who trade more earn lower net returns, and in a separate study \cite{b15} they find that overconfidence shows up as trading intensity in a comparison by gender. Group-level patterns of this kind cannot serve as an individual diagnosis, so Folium avoids demographic shortcuts. Grinblatt and Keloharju \cite{b16} find that investor traits, past performance and the trading situation all affect behavior together, which also limits Folium: one questionnaire and one snapshot of the market can never capture everything.

\section{Machine Learning for Financial Classification}
Friedman \cite{b26} describes gradient boosting as the stage-wise optimization of a function. Chen and Guestrin \cite{b5} introduce XGBoost, a gradient-boosted tree framework with regularization, sparsity-aware split finding and scalable approximate learning. Its strength on tabular financial data is the ability to capture non-linear relationships between unrealized return, market volatility, liquidity and investment horizon without assuming a linear decision boundary. Random forests, introduced by Breiman \cite{b25} as ensembles of randomized decision trees, are a strong non-linear benchmark against which XGBoost can be compared on the same data splits, and the literature does not require boosted trees to outperform forests on artificial behavioral labels. Folium applies this family of models to a multi-class problem in which label quality and probability calibration matter as much as accuracy. Winning against synthetic rules may only show that the model has recovered the labelling process.

Zadrozny and Elkan \cite{b6} study how classifier scores can be turned into class probabilities. This matters because when a model shown to a user expresses certainty, accuracy alone cannot tell whether the stated probabilities match how often it is right. Niculescu-Mizil and Caruana \cite{b27} show that discrimination and probability estimation are separate problems, and that boosted trees in particular push probabilities away from 0 and 1. Guo et al. \cite{b28} show that high accuracy does not imply good calibration and introduce the expected calibration error; their experiments are mostly on neural networks, so no calibration method can be assumed optimal for a tree ensemble. Folium reports calibration error and reliability on held-out data. Calibration only checks agreement with the reference labels; it does not turn an artificial behavioral label into reality.

Saito and Rehmsmeier \cite{b29} show that precision-recall analysis is more informative than ROC analysis under class imbalance. Folium is multi-class, but the same sensitivity to prevalence supports one-vs-rest precision and recall with macro-averaging, since relying on accuracy alone can hide problems with the minority class. López de Prado \cite{b24} discusses evaluation practices for financial machine learning, including the temporal nature of data, leakage and the problems of backtesting. Random splits of artificial data can overestimate generalization, so Folium also uses a temporal holdout.

\section{Explainable Machine Learning}
In the SHAP approach of Lundberg and Lee \cite{b7}, a prediction is attributed to the features relative to a background distribution, and Lundberg et al. later gave an exact and fast algorithm for tree ensembles \cite{b35}. In Folium, SHAP shows whether unrealized loss, crowd or liquidity factors drove the behavioral reading. Its limitation is that SHAP explains the prediction function, not what would happen if a feature were changed in the real world, and the results depend on correlations between inputs. Janzing et al. \cite{b31} frame feature relevance as a causal problem and stress the difference between an observational contribution and an interventional effect. Folium therefore presents SHAP only as predictive attribution, while causal claims require a graph.

Ribeiro et al. \cite{b8} propose LIME, which approximates a black-box model locally with a simpler surrogate. Its faithfulness depends on how perturbations are sampled and how the neighborhood is built. Folium uses tree-specific attributions and so avoids fitting a local surrogate, but its explanations must still be checked numerically and linguistically. Counterfactual explanations by Wachter, Mittelstadt and Russell \cite{b9} specify input changes that would change a prediction, and Mothilal, Sharma and Tan \cite{b30} generate diverse counterfactuals. A change in a prediction is not necessarily an actionable change or a causal effect. Folium's HOLD versus SELL analysis is different: it compares two scenarios on paired market paths and answers a model-specific ``what if'' question, not a full causal question about the investor's real future. Weber, Carl and Hinz \cite{b34} review explainable AI in finance and stress evaluation that is tailored to the financial context. Folium adds bilingual, fact-checked explanation as a design criterion, while recognizing that readability and completeness alone do not prove that explanations improve decisions.

\section{Causal Modeling and Identification}
Pearl \cite{b10} develops structural causal models with directed acyclic graphs, interventions and adjustment criteria. The key lesson is that an observed association is not identified as a causal effect without assumptions about confounding and the graph. Folium models proposed relations between the market state, the portfolio, the investor's constraints and the decision with a directed graph, and its estimates depend on that graph and on the observed covariates; they should never be called causal truth. Robins introduced g-computation \cite{b36}, and Hernán and Robins \cite{b37} give a modern treatment of the identification assumptions of consistency, positivity and conditional exchangeability.

DoWhy, by Sharma and Kiciman \cite{b11}, organizes causal analysis into four stages: modeling the graph, identifying the estimand, estimating, and refuting. Placebo tests, random common causes and subsampling can reveal instability, but passing them does not rule out every unobserved confounder, which is especially relevant for synthetic investor data. Chernozhukov et al. \cite{b32} provide double machine learning for treatment effects with flexible nuisance models; identification still depends on conditional exchangeability, so Folium's synthetic causal checks are not a replacement for a proper identification design.

\section[Long-Horizon Simulation and Portfolio Choice]{Long-Horizon Simulation and\\ Portfolio Choice}
Künsch \cite{b12} develops resampling methods for dependent data, because independent bootstrap samples destroy temporal dependence. This is the reason for sampling contiguous blocks of returns when building market scenarios: block sampling better preserves volatility clustering and serial dependence than drawing independent months, although its accuracy still depends on the historical window, stationarity, the block length and how similar future regimes are to past ones. Politis and Romano's stationary bootstrap \cite{b33} uses random block lengths and is an alternative to fixed blocks. Efron and Tibshirani \cite{b13} formalize the bootstrap for measuring uncertainty, which supports reporting quantiles instead of a single deterministic figure. Folium reports the lower tail, median and upper tail of simulated wealth; these percentiles describe the simulator, not the coverage probabilities of future markets. Hamilton's regime-switching model \cite{b23} reminds us that returns are not drawn from one stable distribution, so sensitivity analysis must allow for changes of regime.

Markowitz \cite{b21} formulates portfolio selection through expected returns, variance and diversification, which suggests that a choice should be judged against the risk of the whole portfolio and not one investment alone. Merton \cite{b22} develops lifetime portfolio selection under uncertainty and shows the importance of horizon, consumption and changing wealth constraints. Folium takes these into account through scenario simulation, but it does not solve the optimal control problem.

\section{Security Foundations for a Financial Web Application}
Folium handles sensitive financial data, so the application layer is built on established standards. JSON Web Tokens are defined in RFC 7519 \cite{b38}, and two-step verification uses time-based one-time passwords as in RFC 6238 \cite{b39}. Passwords are hashed with bcrypt \cite{b40} and their strength is estimated with zxcvbn \cite{b41}. Cross-site request forgery is mitigated with the double-submit cookie pattern described by OWASP \cite{b42}, and per-user isolation uses PostgreSQL row-level security policies \cite{b43}. The tamper-evident audit log follows the hash-chaining idea of Haber and Stornetta \cite{b44} with SHA-256 \cite{b45}, and live progress uses server-sent events \cite{b46}.

\section{Comparison of Existing Approaches}
Table \ref{tab:compare} compares the main kinds of tools available to a retail investor with the approach taken in this project.

\begin{table}[H]
\centering
\small
\begin{tabularx}{\textwidth}{|P{2.9cm}|P{2.1cm}|P{2.3cm}|L|}
\hline
\textbf{Approach} & \textbf{Looks at the decision?} & \textbf{Uses investor's finances?} & \textbf{Weakness for this problem} \\
\hline
Brokerage and portfolio trackers & No & Holdings only & Describe the asset, not whether the action is sensible \\
\hline
Robo-advisers & At onboarding & Risk questionnaire & Do not evaluate an individual pending trade \\
\hline
Behavioral questionnaires & No & Self-reported & Measure a trait, not the present situation \\
\hline
General chat assistants & Yes & Only if typed in & Numbers can be invented; no calibration or audit \\
\hline
Academic bias studies \cite{b2,b14} & After the fact & Account data & Population-level; not a decision-time tool \\
\hline
Folium & Yes, at decision time & Income, outflows, fund, portfolio, horizon & Behavioral labels are synthetic (Section \ref{sec:boundaries}) \\
\hline
\end{tabularx}
\caption{Comparison of existing approaches with Folium}
\label{tab:compare}
\end{table}

\section{Research Gaps}
The reviewed literature in behavioral finance, portfolio theory, calibrated machine learning, causal identification, time-series resampling and explainable AI offers different but complementary methods. Five gaps emerge from it. First, tools and studies tend to treat any loss-making sale or early sale of a winner as a bias, whereas a decision-time tool has to recognize forced sales, rebalancing and near goals as rational. Second, classifiers shown to end users rarely report whether their confidence can be trusted. Third, feature attributions are often presented as causes. Fourth, simple ``what if'' calculators compare actions on different random scenarios or ignore tax, which adds noise and overstates the case for holding. Fifth, generated explanations may misstate figures, and checks that reject wording containing numbers absent from the computed facts are rare.

Behavioral patterns can only be used at the moment of decision if labelling is representative and exceptions are considered. Gradient boosting allows non-linear classification, but reliability and construct validity must be handled separately. Feature attribution improves interpretability without giving causal identification, defensible causal analysis needs graphs and sensitivity analysis, long-horizon resampling must account for temporal dependence and shifts in distribution, and a fluent explanation is not necessarily faithful to the simulation behind it. Folium combines candidate behavior identification, explicit bounds, paired decision comparison and fact-backed explanation in a single decision-time process. This combination is the technical and evaluative contribution of the work; none of the individual algorithms is claimed to be newly invented.
